\documentclass[10pt,a4paper]{article}
\usepackage[margin=1.8cm]{geometry}
\usepackage[T1]{fontenc}
\usepackage{lmodern}
\usepackage{booktabs,tabularx,array}
\usepackage{enumitem}
\usepackage[table]{xcolor}
\usepackage{hyperref}
\usepackage{titlesec}
\usepackage{microtype}
\usepackage{pgfplots}
\usepackage{listings}
\pgfplotsset{compat=1.17}

\setlist{nosep,leftmargin=1.2em}
\titlespacing*{\section}{0pt}{8pt}{3pt}
\titlespacing*{\subsection}{0pt}{6pt}{2pt}
\titleformat{\section}{\large\bfseries}{\thesection}{0.6em}{}
\titleformat{\subsection}{\normalsize\bfseries}{\thesubsection}{0.6em}{}
\setlength{\parskip}{3pt}
\setlength{\parindent}{0pt}
\setlength{\emergencystretch}{2em}
\hypersetup{colorlinks=true,urlcolor=blue,linkcolor=black}
\newcommand{\cf}{ContractForge}
\newcommand{\code}[1]{\texttt{\small #1}}

\definecolor{okgreen}{HTML}{2E7D32}
\definecolor{weakamber}{HTML}{C17A00}
\definecolor{failred}{HTML}{B42318}
\definecolor{rowgray}{HTML}{F2F4F7}

\lstset{basicstyle=\ttfamily\footnotesize,backgroundcolor=\color{rowgray},frame=none,
        columns=fullflexible,keepspaces=true,xleftmargin=4pt,xrightmargin=4pt,aboveskip=4pt,belowskip=4pt}

% one issue = what happens / evidence / why it matters / proposed fix
\newcommand{\what}{\par\smallskip\textbf{What happens.}~}
\newcommand{\evidence}{\par\smallskip\textbf{Evidence.}~}
\newcommand{\why}{\par\smallskip\textbf{Why it matters.}~}
\newcommand{\fix}{\par\smallskip\textbf{Proposed fix.}~}

\title{\vspace{-1.3cm}\textbf{\cf: The L0.5 Feedback Loop on gamma~j}\\[2pt]
\large L0 $\rightarrow$ L0.5 loop $\rightarrow$ L1 $\rightarrow$ L2, and what the loop changes downstream}
\author{Bavya Sri Kilari (2025701029), Akshith K (2024111004), Vivek R (2019112020)\\
\small Autonomous Software Engineering, Monsoon 2026, IIIT Hyderabad}
\date{\small 13 November 2026 \quad Repository: \href{https://github.com/Ohshit-ak/AUTOSE-Leaniers}{\nolinkurl{github.com/Ohshit-ak/AUTOSE-Leaniers}}}

\begin{document}
\maketitle
\vspace{-0.9cm}

% =====================================================================
\section{Summary}
Until now, an L0 spec that failed the L0.5 review was either dropped or regenerated once without being told why
(\code{--on-l0-fail exclude|redo}). The new mode \code{--on-l0-fail loop} regenerates each failed spec with the L0.5
reasons and the rejected spec in the prompt, for up to two rounds. A new spec replaces the old one only if its L0.5 status
improves. Whatever still fails is excluded before L1. We ran it on gamma~j only, with L0 run once and then three L1+L2
repetitions.
\begin{itemize}
\item \textbf{The loop works at L0.5.} It fixed \textbf{6 of 9} failed specs (5 in round 1, 1 in round 2): 5 now pass and 1
      is flagged. Failures drop from 9 to 3, and \textbf{5 more specs} enter L1 (57 $\rightarrow$ 62). It took 8.5 minutes,
      compared with 27 minutes for L0.
\item \textbf{It adds about one verified spec.} Only 1 of the 5 repaired specs is verified end to end. Of the other four,
      two fail PO3 in L2 and two fail PO2 in L1, in all three repetitions (Issue~2).
\item \textbf{The end-to-end count dropped from 52.7 to 47.3, but the loop did not cause it.} L0 was regenerated with a
      changed L0 prompt, and the new L0 was already weaker before the loop ran (57 pass vs 63). On the 17 specs whose text is
      the same in both runs, the results are identical (Issue~1).
\item \textbf{Flagged specs never reach L1}, so turning a fail into a flag has no effect on the results (Issue~3).
\item \textbf{The 3 remaining failures are compile errors}, and round~2 repeats round~1's feedback word for word (Issue~4).
\end{itemize}

\subsection{Run setup}
\begin{tabularx}{\linewidth}{@{}lX@{}}
\toprule
Document    & \code{0000 - gamma j} (74 parsed entries) \\
Model       & \code{qwen2.5-7b-instruct} (Q4\_K\_M, llama.cpp), temperature 0.2 (0.7 on restart); open-weight only \\
Harness     & L0 v10.1 + L0.5 loop (2 rounds, \code{CF\_L05\_ROUNDS}); chain-v4 for L1+L2. Commit \code{8d03bd5}, clean tree \\
Job         & Slurm 76346 (\code{runs/qwen-l05loop-gammaj.sbatch}), GTX 1080 Ti \\
Repetitions & L0 + loop: one run (seed 42). L1+L2: three repetitions (seeds 42, 43, 44), strict only (no fallback run) \\
Baseline    & Frozen L0 \code{qwen-v10.1/rep1} (commit \code{b01f013}) and its chain run \code{qwen-chain-v4} \\
Outputs     & \code{runs/0000 - gamma j/qwen-v10.1-l05loop/} (L0, \code{summary.md}) and \code{qwen-chain-v4-l05loop/} (L1+L2) \\
\bottomrule
\end{tabularx}

\subsection{Terms used in this report}
\begin{tabularx}{\linewidth}{@{}lX@{}}
\toprule
L0.5 status & \emph{pass}: accepted, goes to L1. \emph{flag}: suspicious, kept for review but \emph{not} sent to the chain step.
     \emph{fail}: rejected (does not compile, \code{post} accepts a forbidden output, a condition is always true, \dots). \\
L0.5 loop & Regenerate every \emph{fail} with the reasons and the rejected spec appended to the prompt. Keep the new spec only if
     its status is strictly better (fail $<$ flag $<$ pass). Repeat for at most 2 rounds, then exclude what still fails. \\
repaired spec & A spec that failed L0.5 and passes (or is flagged) after the loop. \\
end to end & Lean accepts one file with the spec, the architecture, PO1--PO3 and \code{theorem system\_ok : Spec.Satisfies \dots}. \\
weak spec & A spec that an implementation made only of copied input fields and constants satisfies. \\
\bottomrule
\end{tabularx}

% =====================================================================
\section{What the loop did at L0.5}
\begin{center}
\begin{tabular}{@{}lrrr@{}}
\toprule
L0.5 status & frozen L0 (\code{b01f013}) & this run, before loop & this run, after loop \\
\midrule
pass                       & 63 & 57 & \textbf{62} \\
flag                       & 3  & 8  & 9  \\
fail                       & 8  & 9  & \textbf{3} \\
\midrule
total                      & 74 & 74 & 74 \\
into L1 (pass only)        & 61\textsuperscript{*} & 57 & \textbf{62} \\
\bottomrule
\end{tabular}
\end{center}
{\footnotesize \textsuperscript{*}The chain step takes specs that pass L0.5 \emph{and} were checked. In the frozen run, 2 of the 63
passing specs (\code{3.4.3/S2}, \code{UC-administrator-manage-plug-in-options}) were not checked, so 61 went to L1.}

\begin{table}[h]
\centering\small
\setlength{\tabcolsep}{4pt}
\rowcolors{2}{white}{rowgray}
\begin{tabularx}{\linewidth}{@{}llcllX@{}}
\toprule
Entry & kind & rounds & after loop & downstream (3 reps) & L0.5 reason fed back \\
\midrule
UC-sales-person-update-product-attributes & use case & 1 & \textcolor{okgreen}{pass} & \textcolor{okgreen}{verified 3/3}
   & four conditions $x \geq 0$ on \code{Nat} fields are always true \\
3.1.3/S2 & requirement & 1 & \textcolor{okgreen}{pass} & \textcolor{failred}{L2 PO3 fails}  & \code{post} accepts a forbidden output \\
5.3/R3   & requirement & 1 & \textcolor{okgreen}{pass} & \textcolor{failred}{L2 PO3 fails}  & \code{post} accepts every changed output \\
5.3/R4   & requirement & 1 & \textcolor{okgreen}{pass} & \textcolor{failred}{L1 PO2 fails}  & \code{post} accepts a forbidden output \\
5.5/R1   & requirement & 1 & \textcolor{okgreen}{pass} & \textcolor{failred}{L1 PO2 fails}  & does not compile (instance synthesis error) \\
UC-administrator-login-administrator & use case & 2 & \textcolor{weakamber}{flag} & not sent to L1 & \code{post} accepts a forbidden output \\
DD-shipping-date               & data   & 2 & \textcolor{failred}{fail} & excluded & does not compile \\
UC-customer-add-item-to-cart   & use case & 2 & \textcolor{failred}{fail} & excluded & does not compile \\
UC-customer-checkout-an-order  & use case & 2 & \textcolor{failred}{fail} & excluded & does not compile \\
\bottomrule
\end{tabularx}
\caption{The 9 entries that failed L0.5 before the loop. \emph{rounds}: loop rounds spent on the entry (an entry fixed in
round~1 is not retried).}
\label{tab:entries}
\end{table}

The loop fixes the failures that a sentence can explain: an always-true condition, or a \code{post} that is too loose (4 of
the 5 passes). It does not fix compile errors: 3 of the 4 compile failures are still failing after two rounds.

% =====================================================================
\section{End-to-end results}
Means over three L1+L2 repetitions (strict). The baseline uses the frozen L0; this run uses the new L0 after the loop.

\begin{center}
\begin{tabular}{@{}lrr@{}}
\toprule
 & frozen L0 $\rightarrow$ \code{qwen-chain-v4} & this run $\rightarrow$ \code{qwen-chain-v4-l05loop} \\
\midrule
specs into L1                         & 61                      & 62 \\
L1 accepted                           & 58.0 (58--58)           & 56.3 (56--57) \\
\quad two-component architectures     & 1.0                     & 0.3 \\
L1 failures: po2 / identity / type    & 0.7 / 2.0 / 0.3         & 4.7 / 1.0 / 0 \\
L2 failures: po3                      & 5.3                     & 9.3 \\
\midrule
\textbf{end to end}                   & \textbf{52.7} (52--53)  & \textbf{47.3} (47--48) \\
\quad on a non-trivial spec           & 30.7 (30--31)           & 26.0 (26--26) \\
\quad requirement / data / use case   & 21.0 / 17 / 14.7        & 20.3 / 17 / 10.0 \\
weak specs                            & 24                      & 23 \\
minutes per repetition                & 11.5                    & 14.0 \\
\bottomrule
\end{tabular}
\end{center}

Within this run, the loop's own contribution is $+1$: without it the 5 repaired specs would have been excluded, and only
\code{UC-sales-person-update-product-attributes} is verified (in all three repetitions). So end to end would have been
$47.3 - 1 = 46.3$ without the loop, and is 47.3 with it. Section~\ref{sec:issues} explains why the drop from 52.7 is not
the loop's doing.

% =====================================================================
\clearpage
\section{Issues, explained}\label{sec:issues}
Each issue gives what happens, the evidence from this run, why it matters, and a proposed fix. None of the fixes have been
made yet.

% ---------------------------------------------------------------------
\subsection*{Issue 1. This run cannot be compared with the frozen run}
\what The job regenerated L0 instead of reusing the frozen \code{qwen-v10.1/rep1}. Both runs use seed 42, but the L0 prompt
has changed since the frozen run: \code{prompt\_hash} is \code{a4ace6f72b38} at \code{b01f013} and \code{08c8321aa2dd} at
\code{8d03bd5}. So two things changed at once, L0 and the loop.

\evidence
\begin{center}
\begin{tabular}{@{}lrrr@{}}
\toprule
Specs in both chain runs (56) & specs & e2e, frozen & e2e, this run \\
\midrule
identical Lean spec text          & 17 & 17.0 & 17.0 \\
different Lean spec text          & 39 & 32.0 & 27.3 \\
\bottomrule
\end{tabular}
\end{center}
Where the spec text is the same, the results are the same. All of the drop comes from specs whose text changed:
\code{1.4/S3}, \code{3.2.3/S1}, \code{5.3/R3}, \code{5.3/R4} and three administrator use cases (install patch, install
plug-ins, remove plug-ins) are lost. \code{3.1.3/S3}, \code{SI-1.2} and \code{UC-administrator-add-user} are gained.
Before the loop ran, the new L0 already had 6 fewer passing specs (57 vs 63).

\why A reader comparing the two headline numbers would conclude that the loop costs 5 verified specs. It adds one. The
difference between the two L0 runs is larger than the effect being measured.

\fix Apply the loop to the \emph{frozen} L0: copy \code{qwen-v10.1/rep1} and run only the check stage with
\code{--on-l0-fail loop}. Then L0 is identical and the loop is the only change. Separately, rerun L0 with the current prompt
over three seeds to see how much L0 varies on its own.

% ---------------------------------------------------------------------
\subsection*{Issue 2. Passing L0.5 after feedback does not mean the spec is implementable}
\what The loop makes a spec acceptable to the L0.5 review, but the repaired specs fail much more often later in the
pipeline than the specs that passed on the first try.

\evidence Repaired specs: 1 of 5 verified (20\%). Specs that passed without the loop: 46.3 of 57 (81\%). The four
failures are the same in all three repetitions:
\begin{itemize}
\item \code{3.1.3/S2}, \code{5.3/R3}: routed single, L2 fails PO3 after 5 attempts.
\item \code{5.3/R4}, \code{5.5/R1}: their new \code{post} relates two output fields, so they are routed to a split, and L1
      fails PO2 in every repetition. The repaired \code{post} relates two output fields, which the router reads as an
      intermediate result.
\end{itemize}
In the frozen run, \code{5.3/R3} and \code{5.3/R4} were routed single and verified with L2 on the first attempt. Their
repaired versions here fail. Those frozen specs were never rejected by L0.5, so the two runs do not show whether
the feedback made the specs harder or the new L0 sample did.
PO2 failures went from 0.7 to 4.7 per repetition, and these two specs account for 2 of them.

\why The loop raises the L0.5 pass count by 5, but that overstates its value: 4 of those 5 specs reach L1/L2 and still fail.

\fix Report the loop's effect at the end-to-end level, not at L0.5. Look at the regenerated \code{post}s of \code{5.3/R4}
and \code{5.5/R1}: if the feedback (``\code{post} accepts a forbidden output'') makes the model add an output-to-output
relation, the split router will send the spec to L1 every time.

% ---------------------------------------------------------------------
\subsection*{Issue 3. Flagged specs never reach L1}
\what The chain step reads only specs with status \emph{pass}. \code{UC-administrator-login-administrator} went from fail to
flag in round~2, and the loop counted it as improved, but it is not used.

\evidence 71 specs are pass or flag after the loop, but only 62 go to L1. All 9 flagged specs are left out, including the
repaired one.

\why Counting fail $\rightarrow$ flag as an improvement makes the loop look more useful than it is in the results.

\fix Either loop on flags too (with the flag reason as feedback), or count only fail $\rightarrow$ pass as improved. Report
the number of flagged specs next to the chain input.

% ---------------------------------------------------------------------
\subsection*{Issue 4. Compile errors are not fixed, and round 2 repeats round 1}
\what When a regenerated spec is not better, the loop keeps the \emph{old} record. Round~2 then builds its feedback from that
old record, so the model gets the same reasons and the same rejected spec as in round~1. It never sees why its round-1
attempt failed.

\evidence \code{DD-shipping-date}, \code{UC-customer-add-item-to-cart} and \code{UC-customer-checkout-an-order} fail to
compile in both rounds, each after 5 attempts per round. Round~2 fixed 1 of 4 entries, against 5 of 9 in round~1.

\why A second round with the same feedback is close to a blind retry. More rounds will cost time without fixing more.

\fix Build round-$k$ feedback from the most recent attempt (its reasons and Lean error), while still keeping the better spec
on disk. For compile errors, include the Lean error at its line, not only ``does not compile''.

% =====================================================================
\section{Suggested order of work}
\begin{enumerate}
\item \textbf{Measure the loop properly} (Issue 1): run the check stage with \code{--on-l0-fail loop} on a copy of the frozen L0,
      then chain-v4 $\times$3. This is a short job and gives the number to report.
\item \textbf{Feedback from the latest attempt} (Issue 4): a small change in \code{\_l05\_loop}, likely to help round~2.
\item \textbf{Decide what to do with flags} (Issue 3): loop on them or stop counting fail $\rightarrow$ flag as improved.
\item \textbf{Then run on the other 17 documents}: the loop is general (it uses only the L0.5 reasons), but report end to end
      with and without it, from the same L0.
\end{enumerate}

\vfill
{\footnotesize Numbers from \code{runs/0000 - gamma j/qwen-v10.1-l05loop/summary.md},
\code{runs/0000 - gamma j/qwen-chain-v4-l05loop/rep\{1,2,3\}/l1l2.json}, \code{runs/0000 - gamma j/qwen-chain-v4/summary.md}
and \code{runs/qwen-l05loop-gammaj.log}.}

\end{document}
